\documentclass{beamer}
\usepackage{amsmath}
\setbeameroption{show notes on second screen=right}
\setbeamertemplate{note page}[plain]

\begin{document}

\begin{frame}
\frametitle{}


$6\frac15 \div 2\frac34$. \note[item]{\alert<1>{Find six and one fifth divided by two and three fourths.}}

\vskip4pt

\pause $= \frac{31}{5} \div \frac{11}{4}$ 
\note[item]{\alert<2>{We rewrite six and one fifth. For the numerator, we have 6 times 5, plus 1. For the denominator, copy the denominator 5. For the second fraction, we rewrite two and three fourths. For the numerator, we have 2 times 4, plus 3. For the denominator, copy the denominator 4.}}

\vskip4pt

\pause $= \frac{31}{5} \times \frac{4}{11}$
\note[item]{\alert<3>{Instead of dividing by eleven fourths, we will multiply by the reciprocal, which is four elevenths.}}

\vskip4pt

\pause $=\frac{31 \times 4}{5 \times 11}$
\note[item]{\alert<4>{When multiplying fractions, we multiply the numerators together and the denominators together. So the numerator is 31 times 4, and the denominator is 5 times 11.}}

\vskip4pt

\pause $=\frac{124}{55}$
\note[item]{\alert<5>{Simplifying the numerator and denominator separately, we have one hundred twenty-four over fifty-five. Since 124 and 55 do not share any common factors, this is our final answer as an improper fraction.}}

\vskip4pt
\pause $= 2\frac{14}{55}$
\note[item]{\alert<6>{We often leave a final answer as an improper fraction, but if needed, we can rewrite it as a mixed number. Since 55 goes into 124 two times, we can rewrite this as a mixed number. The whole number is 2, and the fraction is the remainder, which is 14, over the original denominator, which is 55. Both answers are correct: the improper fraction 124 over 55 is correct, but also, the mixed fraction 2 and 14 over 55 is correct.}}

\end{frame}



%% Blank frames at the end to prevent weird shifts.
\begin{frame}
\end{frame}
\begin{frame}
\end{frame}

\end{document}
